\chapter*{Conclusion G\'{e}n\'{e}rale}
\label{chapter:Conclusion}
\addcontentsline{toc}{chapter}{Conclusion G\'{e}n\'{e}rale}
\setcounter{section}{0}

\section{Bilan g\'{e}n\'{e}ral}

Au terme de ce projet consacr\'{e} \`{a} l'apport de l'Entity Resolution et du Machine Learning dans l'am\'{e}lioration de la qualit\'{e} des donn\'{e}es des op\'{e}rateurs de change \`{a} l'Office des Changes, plusieurs enseignements fondamentaux ont pu \^{e}tre d\'{e}gag\'{e}s.

Ce rapport a permis, dans un premier temps, de pr\'{e}senter le cadre institutionnel de l'Office des Changes, dont la strat\'{e}gie 2025--2029 place la digitalisation et le contr\^{o}le intelligent parmi ses priorit\'{e}s. Dans un second temps, la m\'{e}thodologie du linkage probabiliste multi-champs a \'{e}t\'{e} expos\'{e}e et justifi\'{e}e, et chaque \'{e}tape du pipeline a \'{e}t\'{e} d\'{e}crite en d\'{e}tail, depuis la pr\'{e}paration des donn\'{e}es jusqu'\`{a} la mat\'{e}rialisation des entit\'{e}s. Dans un troisi\`{e}me temps, l'architecture de la plateforme web a \'{e}t\'{e} pr\'{e}sent\'{e}e, avec la communication entre ses composants et ses fonctionnalit\'{e}s de configuration, de suivi, de visualisation et de recherche. Enfin, l'\'{e}valuation exp\'{e}rimentale sur donn\'{e}es synth\'{e}tiques avec v\'{e}rit\'{e} terrain a valid\'{e} l'efficacit\'{e} de l'ensemble.

\section{R\'{e}ponse \`{a} la probl\'{e}matique}

Le syst\`{e}me d\'{e}velopp\'{e} r\'{e}pond \`{a} la probl\'{e}matique pos\'{e}e \`{a} travers cinq contributions concr\`{e}tes. La premi\`{e}re est un pipeline d'Entity Resolution en neuf \'{e}tapes, o\`{u} chaque \'{e}tape est ind\'{e}pendante, configurable et observable, depuis le chargement des d\'{e}clarations jusqu'\`{a} l'\'{e}criture des entit\'{e}s r\'{e}solues. La deuxi\`{e}me est une pr\'{e}paration des donn\'{e}es en deux couches, qui combine un nettoyage g\'{e}n\'{e}rique et une normalisation m\'{e}tier tenant compte des sp\'{e}cificit\'{e}s marocaines des identifiants. La troisi\`{e}me est un moteur probabiliste fond\'{e} sur Splink avec entra\^{i}nement non supervis\'{e}, enrichi de mani\`{e}re optionnelle par la similarit\'{e} s\'{e}mantique neuronale et compl\'{e}t\'{e} par un clustering sur graphe avec mat\'{e}rialisation d'attributs canoniques. La quatri\`{e}me est une plateforme web compl\`{e}te qui rend l'ensemble accessible \`{a} des utilisateurs non informaticiens et qui propose une recherche intelligente d'op\'{e}rateurs. La cinqui\`{e}me est une \'{e}valuation rigoureuse sur jeu synth\'{e}tique avec v\'{e}rit\'{e} terrain, qui \'{e}tablit un F-score de 0,993 et qui \'{e}claire le r\^{o}le de chaque composant par des analyses d'ablation et une validation manuelle.

\section{D\'{e}marche m\'{e}thodologique}

La d\'{e}marche m\'{e}thodologique constitue l'apport le plus durable de ce projet. Elle a montr\'{e} que les approches fond\'{e}es sur un identifiant unique ou sur des comparaisons exactes sont insuffisantes face aux variations et aux valeurs manquantes des donn\'{e}es r\'{e}elles. Elle a \'{e}tabli que la combinaison d'indices faibles par un mod\`{e}le probabiliste permet d'atteindre une qualit\'{e} quasi parfaite sans aucune donn\'{e}e annot\'{e}e. Elle a confirm\'{e} que la pr\'{e}paration des donn\'{e}es et le blocking sont au moins aussi d\'{e}terminants que le mod\`{e}le de scoring lui-m\^{e}me. Et elle a circonscrit avec pr\'{e}cision les situations qui justifient le maintien d'un contr\^{o}le humain, plut\^{o}t que de pr\'{e}tendre \`{a} une automatisation totale.

\section{Limites et perspectives}

L'absence de validation sur donn\'{e}es r\'{e}elles constitue la principale limite actuelle : seule une revue manuelle d'un \'{e}chantillon r\'{e}el par les experts m\'{e}tier permettra de confirmer les performances mesur\'{e}es sur donn\'{e}es synth\'{e}tiques. Le passage \`{a} l'\'{e}chelle vers des volumes de plusieurs millions de d\'{e}clarations devra \^{e}tre test\'{e} et optimis\'{e}. L'hypoth\`{e}se d'ind\'{e}pendance des champs, bien que sans cons\'{e}quence empirique majeure, invite \`{a} la prudence dans l'interpr\'{e}tation des probabilit\'{e}s individuelles.

Les perspectives s'ordonnent naturellement : validation sur donn\'{e}es r\'{e}elles \`{a} court terme, int\'{e}gration dans le processus p\'{e}riodique de chargement et tableau de bord de qualit\'{e} \`{a} moyen terme, extension \`{a} d'autres tables et rapprochement avec des r\'{e}f\'{e}rentiels externes \`{a} long terme.

\bigskip
\noindent En d\'{e}finitive, ce projet aura permis de d\'{e}montrer qu'une probl\'{e}matique de qualit\'{e} des donn\'{e}es peut \^{e}tre adress\'{e}e par une approche technique rigoureuse combinant pr\'{e}paration des donn\'{e}es, Entity Resolution probabiliste et clustering non supervis\'{e}, le tout int\'{e}gr\'{e} dans une plateforme op\'{e}rationnelle et valid\'{e} par une \'{e}valuation rigoureuse. La plateforme est fonctionnelle, l'architecture est extensible, la m\'{e}thodologie est reproductible et les r\'{e}sultats sont \'{e}tablis.
